
\documentclass{article}
\usepackage{colortbl}
\usepackage{makecell}
\usepackage{multirow}
\usepackage{supertabular}

\begin{document}

\newcounter{utterance}

\centering \large Interaction Transcript for game `wordle\_withclue', experiment `medium\_frequency\_words\_clue\_no\_critic', episode 18 with reflexion{-}Qwen3.5{-}9B.
\vspace{24pt}

{ \footnotesize  \setcounter{utterance}{1}
\setlength{\tabcolsep}{0pt}
\begin{supertabular}{c@{$\;$}|p{.15\linewidth}@{}p{.15\linewidth}p{.15\linewidth}p{.15\linewidth}p{.15\linewidth}p{.15\linewidth}}
    \# & \multicolumn{2}{c}{Player} && \multicolumn{2}{c}{Game Master} \\
    \hline

    \theutterance \stepcounter{utterance}  
    & & & \multicolumn{4}{p{0.6\linewidth}}{
        \cellcolor[rgb]{0.9,0.9,0.9}{
            \makecell[{{p{\linewidth}}}]{
                \texttt{\tiny{[P1$\langle$GM]}}
                \texttt{You are a language wizard who likes to guess words by using the given rules.} \\
\\ 
\texttt{Welcome to Wordle! You have six attempts to guess the target word, a valid English word of five lowercase letters (a{-}z). Please use the tags "explanation:" and "guess:" to provide a concise explanation for each guess.} \\
\\ 
\texttt{To help you make an informed guess, you will receive a clue for the word, such as} \\
\texttt{clue: snowy white.} \\
\\ 
\texttt{Here is an example guess based on the clue:} \\
\texttt{explanation: In the fairy tail Snow White, the girl is killed because she eats a poisoned apple. And the word apple has 5 letters.} \\
\texttt{guess: apple} \\
\\ 
\texttt{After each guess, your answer will be validated, and you will receive feedback indicating which letters are correct (green), which letters are correct but in the wrong position (yellow), and which letters are incorrect (red). This feedback can be useful in determining which letters to include or exclude in your next guess.} \\
\\ 
\texttt{For example, the feedback for "apple" might be:} \\
\texttt{guess\_feedback: a<yellow> p<yellow> p<green> l<yellow> e<red>} \\
\\ 
\texttt{The explanation should contain details about how both the guess\_feedback and clue work together in tandem to formulate a new guess.} \\
\\ 
\texttt{Do not generate any other text.} \\
\\ 
\texttt{Let's begin with your first guess.} \\
\\ 
\texttt{clue: court activity} \\
            }
        }
    }
     \\ \\

    \theutterance \stepcounter{utterance}  
    & \multicolumn{4}{p{0.6\linewidth}}{
        \cellcolor[rgb]{0.561,0.737,0.561}{
            \makecell[{{p{\linewidth}}}]{
                \texttt{\tiny{[P1$\rangle$GM]}}
                \texttt{explanation: The clue "court activity" suggests a legal proceeding or event. A common five{-}letter word fitting this description is "trial," which directly refers to a court activity.} \\
\texttt{guess: trial} \\
            }
        }
    }
    & & \\ \\

    \theutterance \stepcounter{utterance}  
    & & & \multicolumn{2}{p{0.3\linewidth}}{
        \cellcolor[rgb]{0.9,0.9,0.9}{
            \makecell[{{p{\linewidth}}}]{
                \texttt{\tiny{[GM$|$GM]}}
                \texttt{attempts = 1} \\
\texttt{target\_word = suing} \\
\texttt{target\_word\_clue = court activity} \\
\texttt{guess = trial} \\
\texttt{guess\_feedback = t<red> r<red> i<green> a<red> l<red>} \\
            }
        }
    }
    & & \\ \\

    \theutterance \stepcounter{utterance}  
    & & & \multicolumn{4}{p{0.6\linewidth}}{
        \cellcolor[rgb]{0.9,0.9,0.9}{
            \makecell[{{p{\linewidth}}}]{
                \texttt{\tiny{[P1$\langle$GM]}}
                \texttt{guess\_feedback: t<red> r<red> i<green> a<red> l<red>} \\
\\ 
\texttt{Provide your response only in this format.} \\
\texttt{explanation: details} \\
\texttt{guess: word} \\
            }
        }
    }
     \\ \\

    \theutterance \stepcounter{utterance}  
    & \multicolumn{4}{p{0.6\linewidth}}{
        \cellcolor[rgb]{0.561,0.737,0.561}{
            \makecell[{{p{\linewidth}}}]{
                \texttt{\tiny{[P1$\rangle$GM]}}
                \texttt{explanation: The feedback for "trial" shows only 'i' is correct and in the right position, while 't', 'r', 'a', and 'l' are incorrect. Since the clue is "court activity" and 'i' is confirmed, I need a word related to law or proceedings that contains 'i' as the third letter. "Jury" fits the context of a court activity and has 'u' as the second letter, but the third letter is 'r', not 'i'. "Gavel" is an object, not an activity. "Hear" is too short. "Issue" has 'i' as the second letter. "Indict" is six letters. "Inquiry" is seven. Let's reconsider words with 'i' in the third spot: "sight", "light", "night", "fight", "right", "might", "tight", "pilot", "vital", "civil", "diver", "river", "liver", "giver", "wider", "biter", "cider", "fiver", "hider", "liver", "miler", "piper", "rider", "siver", "titer", "viper", "wiper", "yiver". Among these, "civil" relates to court activity (civil court). "Jury" doesn't fit the 'i' position. "Hear" is 4 letters. "Trial" failed. How about "issue"? No, 'i' is first. "Inquiry"? Too long. "Arbit"? No. "Judge"? No 'i'. "Jailed"? No. "Jailed" has 'i' as third? J{-}A{-}I{-}L{-}E{-}D, 6 letters. "Jail" is 4. "Prison" is 6. "Court" is 5 but 'o' is second. "Hear" is 4. "Test" is 4. "Case" is 4. "Suit" is 4. "Laws" is 4. "Rule" is 4. "Order" is 5, 'r' is second. "Ruling" is 6. "Verdict" is 7. "Sentence" is 8. "Plea" is 4. "Gavel" is 5, 'a' is second. "Gavel" has no 'i'. "Bench" is 5, 'e' is second. "Judge" is 5, 'u' is second. "Jury" is 4. "Witness" is 7. "Lawyer" is 6. "Attorney" is 8. "Prosecutor" is 10. "Defendant" is 9. "Plaintiff" is 9. "Counsel"} \\
            }
        }
    }
    & & \\ \\

    \theutterance \stepcounter{utterance}  
    & & & \multicolumn{2}{p{0.3\linewidth}}{
        \cellcolor[rgb]{0.9,0.9,0.9}{
            \makecell[{{p{\linewidth}}}]{
                \texttt{\tiny{[GM$|$GM]}}
                \texttt{Error: The response should contain the 'guess:' keyword exactly once.} \\
            }
        }
    }
    & & \\ \\

    \theutterance \stepcounter{utterance}  
    & & & \multicolumn{2}{p{0.3\linewidth}}{
        \cellcolor[rgb]{0.9,0.9,0.9}{
            \makecell[{{p{\linewidth}}}]{
                \texttt{\tiny{[GM$|$GM]}}
                \texttt{game\_result = ABORT} \\
            }
        }
    }
    & & \\ \\

    \theutterance \stepcounter{utterance}  
    & & \multicolumn{4}{p{0.6\linewidth}}{
        \cellcolor[rgb]{0.95,0.95,0.95}{
            \makecell[{{p{\linewidth}}}]{
                \texttt{\tiny{[P1$\leftrightarrow$P]}}
                \texttt{Post{-}episode Reflexion (not sent to GM):} \\
\texttt{New plan: The last attempt was aborted because the response contained excessive text, a long list of brainstormed words, and incomplete sentences, violating the strict format requirement to provide only a concise explanation followed immediately by the guess. Additionally, the explanation failed to explicitly integrate the clue "court activity" with the feedback logic. Next time, I will strictly adhere to the format: a single, concise sentence in the "explanation:" tag that references the clue and the specific feedback (e.g., "Since 'i' is the only correct letter in the third position for 'court activity', I will try a word like 'judge' or 'order'"), followed immediately by a valid five{-}letter guess in the "guess:" tag without any extra commentary or lists.} \\
            }
        }
    }
    & \\ \\

\end{supertabular}
}

\end{document}
